\documentclass[12pt]{book}

\input{codigofuente/preambulo}
\usepackage{makeidx}

\renewcommand{\gcd}{\operatorname{mcd}}


%\begin{comment}
%\usepackage{draftwatermark}
%\SetWatermarkText{\textsc{No enviar}} % por defecto Draft 
%\SetWatermarkScale{0.9} % para que cubra toda la página
%\SetWatermarkColor[rgb]{1,0,0} % por defecto gris claro
%\SetWatermarkAngle{55} % respecto a la horizontal
%\end{comment}


\usepackage{blindtext}


%\makeindex
\begin{document}

\mainmatter


\frontmatter
\input{codigofuente/portada}
\renewcommand{\contentsname}{Índice}


\chapter*{Agradecimientos}

\pagenumbering{roman}
\addcontentsline{toc}{chapter}{Agradecimientos }


\newpage

\chapter*{Dedicatoria}

\addcontentsline{toc}{chapter}{ Dedicatoria}





\chapter*{Resumen}

\addcontentsline{toc}{chapter}{Resumen }


d
\newpage 





\tableofcontents
\mainmatter





\chapter{Introducción}

\newpage


\section{Antecedentes}




\newpage




	
%\newpage

%\section{Justificación}


\newpage

\section{Objetivo}


\subsection{ Objetivo General.}



\subsection{ Objetivos especificos.}

\begin{enumerate}
\item [$\bullet$] 
	\end{enumerate}
	
	\newpage
	
	\section{Pregunta de investigación}
		
\chapter{Marco Teórico}


\section{Números enteros y el algoritmo de división}















\chapter{Métodologia}

En este capítulo se presenta la metodologíay naturaleza de la investigación, aplicada para el logro de los objetivos propuestos. Además, comprende el tipo de investigación, tipo de análisis, método y la organización del trabajo.

\section{Tipo de Investigación.}
%Esta investigación es de tipo documental, donde se procederá hacer un análisis de las diferentes fuentes bibliográficas que están relacionadas a la temática, entre ellas podemos mencionarlos libros y artículos científicos. Se pretende una investigación de tipo axiomática, cuyo interés es estudiar los diferentes enunciados o proposiciones, que fundamentarán la base de las hipótesis que serán demostradas en la investigación.


\section{Método de la investigación.}

%Para este estudio se utilizó el método descriptivo, inductivo y comparativo para la obtención de las informaciones. En cuanto al análisis descriptivo se tomaron en cuenta fuentes escritas, como libros, artículos escritos por autores reconocidos que han hecho investigaciones relacionadas con la conjugación primaria en semigrupos o temas relacionados. En cuanto al análisis inductivo se fueron enlazando las informaciones para llegar a conclusiones generales. De igual manera se realizaron comparaciones con informaciones proporcionadas por los diferentes autores ya citados.



\section{Diseño de la investigación.}

%El diseño de la investigación para dar respuestas a las interrogantes planteadas inicialmente vendrá dado por la revisión y análisis de las diferentes fuentes de información, tales como artículos, libros, tesis, trabajos de investigación. Estas fuentes de información son las que corroborarán el conocimiento inmediato de la investigación.




\chapter{Presentación de Resultados}





En este capítulo se presentan los resultados del estudio realizado donde
se da respuestas a los objetivos de la investigación.\\

\chapter{Conclusión y Recomendaciones}

En este capítulo presentamos las conclusiones a la que se arribaron, de igual manera se sugieren algunas líneas de investigación para el desarrollo futuro en el área de Álgebra Abstracta. Finalmente se presenta las referencias bibliográficas, fuentes que sirvieron de base para la recopilación de las informaciones en esta investigación.

\section{Conclusiones}

\section{ Recomendaciones}





%    Insert the bibliography data here.
%\bibliographystyle{amsplain}
%\bibliography{Bibitw.bib}


\begin{thebibliography}{12}

\bibitem[Bell]{Bell} E.T. Bell. Men of Mathematics. (Victor Gollancz Ltd., London; 1937.
\
\end{thebibliography}


%\addcontentsline{toc}{chapter}{Índice alfabético} 


%\printindex

\end{document}